\section{Signed rigidity on simple spaces}
\label{sec:signed}

The meromorphic section furnished by Theorem~\ref{nv:meromorphic} need
not have an effective divisor.  We therefore need a boundary argument
that retains both its zeroes and its poles.  The result below supplies
that argument.

\begin{theorem}[Signed rigidity]\label{signed:torsion}
Let \((X,D)\) be a dlt pair on a normal irreducible compact Kähler
space, and suppose that the pair has the resolution property of
Definition~\ref{def:lc-strata-resolution}.  Suppose also that \(X\) is
globally \(\Q\)-factorial, that \(a(X)=0\), that \(X\) is simple,
and that \(D=\sum_{i=1}^sD_i\) is reduced.  Put \(L=K_X+D\), and
assume the following.
\begin{enumerate}
\item The rational line \(L\) is analytically nef and has an actual
rational linear equivalence
\[
 L\sim_{\Q}G=\sum_{i=1}^s a_iD_i,\qquad a_i\in\Q.
\]
\item For some real \(c>0\), the pullback of the class
\(\{L-cD\}\) to a projective resolution with smooth compact
Kähler source is pseudo-effective.
\item The rational holomorphic line \(L|_D\) is semiample on the
whole reduced analytic space \(D\): some Cartier multiple is
generated at every point of \(D\).
\end{enumerate}
Then \(L\) is torsion.  More precisely, for some positive integer
\(q\), the holomorphic line \(\OO_X(qL)\) is isomorphic to
\(\OO_X\).
\end{theorem}

To apply this theorem in the proof of Proposition~\ref{simple:case},
Theorem~\ref{nv:meromorphic} supplies a signed representative of the
canonical bundle.  After resolving, we enlarge the support of the
transformed boundary and this signed canonical divisor to a reduced SNC
divisor; the resulting adjoint remains pseudo-effective.
Corollary~\ref{prog:signed-nef} gives an ordinary dlt nef model
\((X_{\mathrm{nef}},D_{\mathrm{nef}})\) whose adjoint
\(L_{\mathrm{nef}}=K_{X_{\mathrm{nef}}}+D_{\mathrm{nef}}\) has a
signed representative supported on \(D_{\mathrm{nef}}\), supplying
the first hypothesis.  Theorem~\ref{floor:generation} supplies
semiampleness on the whole reduced \(D_{\mathrm{nef}}\).  On a common
smooth resolution, simplicity excludes uniruledness, so Ou's criterion
\cite[Theorem~1.1]{Ou2025} and exceptional translation for the canonical
comparison give pseudo-effectivity of the pullback of
\(K_{X_{\mathrm{nef}}}=L_{\mathrm{nef}}-D_{\mathrm{nef}}\).  Thus the
second hypothesis holds with \(c=1\).  Once signed rigidity makes the
actual line \(L_{\mathrm{nef}}\) torsion, the program comparison and
Lemma~\ref{bir:unique-current} let us subtract the added boundary and give
the required decomposition for the original adjoint.  The final subsection
verifies these steps.

The signed construction in \cite[Proposition~3.3]{LI} is the source
of the root and separation operations used below.  The lifting
method is that of \cite[Sections~10--12]{BL}.
We will establish the analytic form of the signed construction and
the change to the lifting argument caused by its residual poles.
The numerical argument first singles out positive-dimensional fibers in
the positive boundary.  We then lift those fibers through every finite
boundary neighborhood and deform them to compact subspaces outside the
boundary. An argument using the cycle space and Baire's theorem turns these
deformations into a covering family, which simplicity excludes.

\subsection{The boundary detected by the nef class}

If \(D=0\), the signed equivalence already gives
\(L\sim_{\Q}0\).  We may therefore assume in the constructions
below that \(D\ne0\); in particular \(\dim X>0\).

Fix the data of Theorem~\ref{signed:torsion}, and write
\[
 G=G_+-G_-,\qquad D_0=\sum_{a_i=0}D_i,
\]
where \(G_+\) and \(G_-\) are effective with disjoint prime supports.
Choose a positive integer \(m\) so that \(mG_+\), \(mG_-\),
\(mD_0\), and \(mL\) are Cartier, the given rational equivalence
induces a fixed isomorphism
\[
 N_0:=\OO_X(mG)\simeq\OO_X(mL),
\]
and \(N_0|_D\) is generated.  Its sections give a morphism and an
actual line isomorphism
\begin{equation}\label{signed:boundary-map}
 \phi:D\longrightarrow P=\PP^b,\qquad
 N_0|_D\simeq\phi^*\OO_P(1).
\end{equation}
Let \(n=\dim X\).  Choose a projective resolution
\(\mu:\widetilde X\to X\) as in the second hypothesis of the theorem,
and a Kähler class \(\omega\) on \(\widetilde X\).  The numerical
dimension of the nef pullback is
\[
 v=\nu(L):=\max\bigl\{k\in\{0,\ldots,n\}:
       (\mu^*\{L\})^k\omega^{n-k}>0\bigr\}.
\]

The next argument is the Kähler form of
\cite[Lemma~3.2]{LI}.

\begin{lemma}\label{signed:boundary-geometry}
If \(v=0\), then \(D=0\) and \(L\sim_{\Q}0\).  If \(v=n>0\),
then \(a(X)=n\).  If \(0<v<n\) and \(r=v-1\), then
\[
 \dim\phi(D_i)\leq r\quad\text{for every }i,
 \qquad
 \max_{a_i>0}\dim\phi(D_i)=r,
\]
and
\begin{equation}\label{signed:small-intersections}
 \dim\phi\bigl(\Supp G_+\cap(\Supp G_-\cup D_0)\bigr)<r.
\end{equation}
For \(r=0\), the intersection in
Equation~\eqref{signed:small-intersections} is empty.
\end{lemma}

\begin{proof}
Choose a Kähler class \(\theta\) on \(X\).  We use Cartier
intersection products downstairs, computed on a resolution by
projection formula.  These products detect \(v\).  Indeed, on the
chosen resolution the class \(h=\mu^*\theta\) is nef and big.
Choose \(C,\epsilon>0\) such that
\(C\omega-h\) is nef and \(h-\epsilon\omega\) is
pseudo-effective.  Pairing these two inequalities successively with
products of the nef classes \(\mu^*\{L\}\), \(h\), and
\(\omega\) shows that
\[
 (\mu^*\{L\})^k h^{n-k}>0
 \quad\Longleftrightarrow\quad
 (\mu^*\{L\})^k\omega^{n-k}>0.
\]
We may therefore use \(\theta\) in all the tests defining \(v\).

Suppose first that \(v<n\).  The pseudo-effective pullback in the
theorem can be paired with nef products, so
\begin{equation}\label{signed:boundary-zero}
 0\leq (\{L\}-c\{D\})\{L\}^{v}\theta^{n-v-1}
   =-c\sum_i\{D_i\}\{L\}^{v}\theta^{n-v-1}.
\end{equation}
Each summand on the right is nonnegative.  For example, pull the
generated line in Equation~\eqref{signed:boundary-map} and
\(\theta\) to a resolution of \(D_i\); their representatives are
semipositive and positive at general smooth points, respectively.
Thus every summand is zero.  When \(v=0\), a nonzero effective
divisor has strictly positive \(\theta^{n-1}\)-degree.  Hence
\(D=0\), and the signed equivalence gives \(L\sim_{\Q}0\).

Now let \(0<v<n\).  For any irreducible effective cycle \(V\) in
\(D\), the mixed degree
\(\{L\}^k\theta^{\dim V-k}\cdot V\) is positive exactly when
\(\dim\phi(V)\geq k\).  To see this, resolve \(V\) and use a
Fubini--Study representative for \(m\{L\}|_V\).  Its generic rank
is \(\dim\phi(V)\), and its wedge with the remaining Kähler
factors is positive on a nonempty open set exactly in the asserted
range.  The vanishing in Equation~\eqref{signed:boundary-zero}
therefore gives \(\dim\phi(D_i)\leq r\).  On the other hand,
\begin{equation}\label{signed:positive-attainment}
 0<\{L\}^{v}\theta^{n-v}
   =\sum_i a_i\{D_i\}\{L\}^{r}\theta^{n-v}.
\end{equation}
All the component tests here are nonnegative, so a component with
\(a_i>0\) attains dimension \(r\).

It remains to separate its intersections from a general fiber.
On \(H^{1,1}_{\BC}(\widetilde X,\R)\) consider the form
\[
 q(\alpha,\beta)=
 \alpha\beta(\mu^*\{L\})^r h^{n-r-2},\qquad
 Q_{ij}=q(\mu^*\{D_i\},\mu^*\{D_j\}).
\]
The mixed Kähler Hodge index theorem, followed by approximation of
the nef factors by Kähler classes, says that \(q\) has at most one
positive direction; see \cite[Theorems~A and C]{DinhNguyen}.  Put
\(l=\mu^*\{L\}\).  The definition of \(v\) and
Equation~\eqref{signed:boundary-zero} give
\[
 q(l,l)=0,\qquad q(l,\mu^*\{D_i\})=0,
 \qquad q(l,h)=\{L\}^{v}\theta^{n-v}>0.
\]
Linear algebra now makes \(q\) negative semidefinite on
\(l^\perp\): a positive vector in \(l^\perp\), together with a
suitable vector in the span of \(l,h\), would give two positive
directions.  In particular \(Q\) is negative semidefinite.

For \(i\ne j\), one has \(Q_{ij}\geq0\).  This assertion uses
the effective intersection cycle of the distinct \(\Q\)-Cartier
divisors \(D_i,D_j\) on \(X\), and projection formula.  It does
not require their total transforms to intersect effectively.  The
two Cartier equations have a proper intersection: the dlt ambient
space is klt, hence Cohen--Macaulay, after dropping the reduced
boundary, and the distinct divisor equations form a regular
sequence.  The intersection cycle is consequently pure of
codimension two, and each of its mixed nef degrees is nonnegative.

Write the coefficient vector as
\(\mathbf a=\mathbf a^+-\mathbf a^-\), with nonnegative vectors
of disjoint supports.  The signed equivalence and orthogonality to
\(l\) give \(Q\mathbf a=0\).  Thus
\[
 Q(\mathbf a^+,\mathbf a^+)
   =Q(\mathbf a^+,\mathbf a^-)\geq0.
\]
Negative semidefiniteness forces equality and
\(Q\mathbf a^+=0\).  For every \(j\) with \(a_j\leq0\), its
row is a sum of nonnegative off-diagonal terms:
\[
 0=\sum_{a_i>0}a_iQ_{ji}.
\]
Each term is zero.  Applying the generated-line degree test to the
effective intersection cycles proves
Equation~\eqref{signed:small-intersections}.  If \(r=0\), any
nonempty intersection would have positive \(\theta^{n-2}\)-degree,
so these intersections are empty.

Finally, if \(v=n>0\), the nef volume criterion makes the rational
line on \(\widetilde X\) big
\cite[Theorem~0.5]{DemaillyPaun2004}.  A big holomorphic line on a
compact Kähler manifold has maximal Iitaka dimension.  Hence
\(a(\widetilde X)=n\), and birational invariance gives \(a(X)=n\).
\end{proof}

Thus, when \(0<v<n\), a positive component attaining image dimension
\(r=v-1\) has general fibers of dimension \((n-1)-r=n-v>0\).
Equation~\eqref{signed:small-intersections} makes those fibers disjoint
from \(\Supp G_-\cup D_0\).  The next construction prepares their normal
directions and adjunction for infinitesimal lifting.

\subsection{Local roots and residue with residual poles}

We work for the rest of the signed argument in the range
\(0<v<n\), with \(r=v-1\).  Choose \(\ell>1\) divisible by \(m\)
and by every nonzero integer \(|ma_i|\), and put
\(a=\ell/m\), a positive integer.

\begin{lemma}[Analytic signed charts]\label{signed:root-charts}
For each \(t\in P\), after shrinking an open neighborhood
\(U\subset P\), choose a line \(R_U\) together with an isomorphism
\(R_U^\ell\simeq\OO_P(1)|_U\).  A comparison of two such roots on an
overlap is power-compatible if its \(\ell\)-th power commutes with these
isomorphisms.  If
\(t\notin\phi(\Supp G_+)\), one may take \(U\) disjoint from this
closed image and set \(Z_U=S=T=\varnothing\); the assertions below are
then vacuous for the unique maps from the empty spaces.  If
\(t\in\phi(\Supp G_+)\), there are a Hausdorff normal
analytic space \(Z_U\) of pure dimension \(n\), considered near a
reduced Cartier divisor \(S\) of pure dimension \(n-1\), a
reduced Weil divisor \(T\) with no component in common
with \(S\), and maps
\[
 \pi:Z_U\longrightarrow X,\qquad g:S\longrightarrow U
\]
with the following properties.
\begin{enumerate}
\item The map \(g\) is proper, and \(\pi|_S\) is finite over
\(D_U=\phi^{-1}(U)\), with \(g=\phi\pi|_S\).  Its image covers
the positive components over \(U\), and
\[
 \OO_S(S)\simeq g^*R_U.
\]
The image of \(S\cap T\) lies in
\(\phi(\Supp G_+\cap(\Supp G_-\cup D_0))\).
Near \(S\), the underlying sets satisfy
\[
 \pi^{-1}(|D|)=|S|\cup|T|.
\]
\item The pair \((Z_U,S+T)\) is log canonical, and there is a fixed
actual isomorphism of divisorial sheaves
\begin{equation}\label{signed:adjoint-chart}
 \tau:\OO_{Z_U}(aS)\xrightarrow{\ \simeq\ }\omega_{Z_U}(S+T).
\end{equation}
The spaces \(Z_U\) and \(S\) are Cohen--Macaulay.  The support
of \(T\) is locally set-theoretically principal.
Off \(T\), the space \(Z_U\) is klt and its canonical sheaf is
invertible.  The map \(\pi\) is quasi-finite near \(S\setminus T\).
\item Projective log resolutions of the total divisor data can be
chosen with Kähler forms on neighborhoods of the compact sets over
each compact fiber of \(g\).  Power-compatible changes of the
boundary root extend to isomorphisms of ambient germs that preserve
the full ideal \(\OO_{Z_U}(-S)\), the reduced divisor \(T\), and
Equation~\eqref{signed:adjoint-chart}.  Resolutions may be chosen
functorially for these germ isomorphisms.
\end{enumerate}
Here the nonempty space \(Z_U\) is a neighborhood of the whole compact
fiber under consideration; no extension of \(g\) to \(Z_U\) is asserted.
\end{lemma}

\begin{proof}
The empty branch follows by shrinking away from the closed image
\(\phi(\Supp G_+)\).  We construct the nonempty charts in four stages:
take roots and trivialize the remaining adjoint torsion, separate the
positive and negative root divisors, retain the root line that controls the
normal direction, and pass to an ordinary neighborhood of the compact fiber.

First consider an analytic open set on which the Cartier lines of
\(mG_+\), \(mG_-\), and \(mD_0\) are trivial, and let
\(f_+,f_-,f_0\) be the functions representing their sections.  Normalize
the entire finite space
\[
 \{y_+^\ell=f_+,\qquad y_-^\ell=f_-,\qquad y_0^\ell=f_0\},
\]
retaining all its components and the lifted action of \(\mu_\ell^3\)
that multiplies the three root coordinates.  On an overlap, the functions
\(f_+,f_-,f_0\) change by holomorphic units.  Local \(\ell\)-th roots of
those units give comparisons of the finite spaces that lift uniquely to
their normalizations; different choices differ by \(\mu_\ell^3\).
The quotient stacks of these normalized charts by \(\mu_\ell^3\) glue to
the simultaneous normalized root stack over \(X\), which we denote by
\(\mathcal X_{\mathrm{rt}}\).  The zero divisors
of the three root sections are denoted by \(S_+,S_-,S_0\).  At a
generic prime of multiplicity \(b\), a normalized chart of
\(y^\ell=x^b\) has ramification index \(\ell/b\) and
\(\ord(y)=1\).  Here \(b=|ma_i|\) or \(b=m\), so it divides
\(\ell\).  Thus each root divisor is generically reduced and
Cartier.  A Cartier divisor on a normal space satisfies Serre's
condition \(S_1\), so it is reduced.  Log ramification with the
full reduced boundary gives a log canonical pair and the rational
equivalence
\[
 K_{\mathcal X_{\mathrm{rt}}}+S_++S_-+S_0
       \sim_{\Q}a(S_+-S_-).
\]
The ambient charts of \(\mathcal X_{\mathrm{rt}}\) are klt.  Off the
boundary this follows from
the klt complement downstairs and the unramified charts.  For a
place centered in the Cartier boundary, dropping that boundary
increases its log discrepancy by its strictly positive order;
this also makes every zero-discrepancy place positive.

To turn the rational adjoint equivalence into an actual isomorphism,
consider the integral reflexive sheaf on \(\mathcal X_{\mathrm{rt}}\)
\[
 \mathcal F=
 \bigl(\omega_{\mathcal X_{\mathrm{rt}}}(S_++S_-+S_0)\otimes
       \OO_{\mathcal X_{\mathrm{rt}}}(-a(S_+-S_-))\bigr)^{**}.
\]
It is torsion.  Choose a periodicity
\(\mathcal F^{[q]}\simeq\OO\), take the relative spectrum of its
reflexive-power algebra, and normalize; write \(\rho_{\mathrm{ind}}\)
for this finite map, and retain \(S_+,S_-,S_0\) for the pulled-back
root divisors on the cover.  At codimension one this
is a cover \(w^q=u\) for a unit \(u\), hence is unramified.
Log canonicity, the klt ambient property, and reduced Cartier
root divisors persist.  The tautological evaluation of this
algebra trivializes the reflexive pullback
\((\rho_{\mathrm{ind}}^*\mathcal F)^{**}\): it does so in
codimension one, and the identity then extends reflexively.
The evaluation is a sheaf morphism on the cover stack itself,
so the resulting actual divisorial isomorphism
\[
 \omega(S_++S_-+S_0)\simeq\OO\bigl(a(S_+-S_-)\bigr)
\]
on the cover is equivariant on every atlas.

Next separate the positive and negative root divisors.  Let
\(\mathcal X_{\mathrm{sep}}\) be the normalization of the blowup of the
ideal generated by their equations on the cover, and write \(p_1\) for its
map.  This blowup embeds in a relative
\(\PP^1\); its normalization still has fibers of dimension at
most one.  Each exceptional prime therefore lies over a
codimension-two positive--negative intersection.  At its generic
point the original dlt pair is a two-branch SNC pair.  After
extracting roots of units the normalized Kummer chart there has
smooth coordinates
\(x=u^{\ell/b_+}\), \(z=w^{\ell/b_-}\), where \(b_+\) and
\(b_-\) are the two corresponding multiplicities.  On this chart
\(\mathcal F\) is a line and its periodicity cover is unramified.
The two-coordinate blowup calculation consequently applies at
every exceptional generic point.  If \(E\) denotes its
exceptional divisor, it gives reduced Cartier divisors
\[
 E,\qquad S'_+=p_1^*S_+-E,\qquad S'_-=p_1^*S_--E,
\]
with \(S'_+\cap S'_-=\varnothing\).  No such two-branch stratum
is contained in \(S_0\), so \(p_1^*S_0\) is reduced and has no
exceptional component.  The full boundary
\(S'_++S'_-+E+p_1^*S_0\) is crepant and reduced.  Its
codimension-one equality with the pullback adjoint gives the
discrepancy equality for every further valuation, hence log
canonicity.  The ambient charts after this blowup are still
klt.  Off the full boundary the blowup is an isomorphism to
the old klt complement.  A place of zero pair discrepancy
centered in that boundary gains its strictly positive order
when the effective Cartier boundary is dropped, while a
place of positive pair discrepancy remains positive.  The
strict divisor \(S'_+\) is finite over \(S_+\):
before normalization it lies in one section of the projective
ratio coordinate, so its proper map has finite fibers, and
normalization is finite.

On \(\mathcal X_{\mathrm{sep}}\), put \(S=S'_+\) and, temporarily,
\(T=S'_-+E+p_1^*S_0\).  Write
\(\pi_{\mathrm{sep}}:\mathcal X_{\mathrm{sep}}\to X\) for the composite
map, and define the retained root line
\[
 \mathcal R:=\OO_{\mathcal X_{\mathrm{sep}}}(S'_+-S'_-),
 \qquad \mathcal R^\ell\simeq\pi_{\mathrm{sep}}^*N_0.
\]
Near \(S\), the negative strict root section is a unit.  The fixed adjoint
isomorphism becomes
\(\omega_{\mathcal X_{\mathrm{sep}}}(S+T)
  \simeq\OO_{\mathcal X_{\mathrm{sep}}}(aS)\), and the positive root
section divided by that negative unit section is a section of
\(\mathcal R\) with zero divisor \(S\).  This line and section determine
the Cartier divisor \(S\) and its normal line; they must survive passage
to an ordinary chart.

The pair consisting of \(\mathcal R\) and its displayed power isomorphism
defines a map to the gerbe \(\sqrt[\ell]{N_0}\), which parametrizes
\(\ell\)-th roots of \(N_0\).  Take the relative coarse space over this
gerbe.  On an atlas this quotients by the finite relative inertia kernel,
the subgroup of each stabilizer acting trivially on \(\mathcal R\).
Thus the quotient retains the root character of \(\mathcal R\).  The
deck transformations of the representable periodicity cover
are not this inertia kernel and are not removed.

Both \(\mathcal R\) and its section descend through the kernel, so the
reduced image \(S\) remains Cartier and its full ideal descends.  We retain
\(S,T\) for the reduced images.  A finite-group norm of
a local equation of the upstairs \(T\) has precisely its image
as zero set.  Thus the reduced image \(T\) is locally
set-theoretically principal; it need not be Cartier.  All
divisorial ramification lies on the full boundary.  Log
ramification identifies the invariant log dualizing sheaf in
codimension one with the downstairs one.  Taking reflexive
hulls and invariants of the equivariant adjoint isomorphism
gives an actual descended isomorphism \(\OO(aS)\simeq\omega(S+T)\).
Its ordinary pullback below is Equation~\eqref{signed:adjoint-chart}.
Outside \(S+T\), the local
finite quotient is quasi-étale, so its klt upstairs complement gives
a klt downstairs complement by finite discrepancy comparison.  Away from
\(T\), dropping the Cartier divisor \(S\) from the lc pair
increases the log discrepancy of every place centered in \(S\) by
its strictly positive order.  Thus the quotient charts are klt away from
\(T\), and the same isomorphism makes their canonical sheaves invertible
there.  Off \(T\) the separating blowup is an
isomorphism, while the other operations are finite; hence
the composite map to \(X\) is quasi-finite near \(S\setminus T\).

The same quotient preserves Cohen--Macaulayness of both the
ambient space and \(S\).  Before the quotient the ambient
space is klt and hence Cohen--Macaulay, and \(S\) is Cartier
there.  Locally its finite quotient ring, and the quotient
ring of \(S\), are invariant direct summands of those finite
Cohen--Macaulay rings.  A system of parameters downstairs is
one upstairs; the vanishing of lower local cohomology upstairs
and the invariant projection give its vanishing for the
downstairs direct summand.  This proves the stated depth.

Removal of the relative inertia kernel also makes the strict
positive divisor finite and representable over
\[
 D\times_P\sqrt[\ell]{\OO_P(1)}.
\]
Its root line is \(\OO_S(S)\), and its image covers
\(\Supp G_+\).  The intersection \(S\cap T\) maps to the
positive--negative or positive--zero intersections: the only
new component is the separating exceptional divisor.  Before the quotient,
the inverse image of \(|D|\) is the support of
\(S'_++S'_-+E+p_1^*S_0\).  The finite quotient sends this full support
to \(|S|\cup|T|\), and the equality of underlying sets persists under the
ordinary base change below.  The image of \(S\cap T\) in \(P\) consequently
has dimension \(<r\), by Lemma~\ref{signed:boundary-geometry}.

Finally we realize these stack data on an ordinary analytic neighborhood
of the whole compact fiber.  We use the dimension-free root-neighborhood
result \cite[Lemma~10.1]{BL}.
For a compact analytic subspace \(F\) of a Hausdorff analytic
space, that lemma extends a specified \(\ell\)-th root of a
line on \(F\) to a root on a neighborhood.  It also extends
a specified power-compatible comparison of two roots uniquely
as a germ about \(F\).  Its proof is the analytic Kummer
sequence and continuity of the cohomology of \(\mu_\ell\)
over neighborhoods of a compact set, in degrees two, one,
and zero.

Apply it to \(F=\phi^{-1}(t)_{\mathrm{red}}\subset X\) with
the root prescribed by \(R_U\) in
Equation~\eqref{signed:boundary-map}.  Apply the comparison
clause also inside \(D\).  Properness of \(\phi\) allows a
shrink of \(U\) for which the comparison is defined along all
of \(D_U\): remove the closed image of its complement.
Pulling the relative coarse construction across this root
over the whole neighborhood of \(F\) gives the ordinary Hausdorff space
\(Z_U\).  Here the chosen root defines a map from that neighborhood to
\(\sqrt[\ell]{N_0}\), and the relative quotient is representable over
the gerbe, so this pullback is an ordinary space.  Finite invariant algebras
glue its local quotient charts.  The preceding finite representable map
becomes the finite map \(S\to D_U\), so \(g\) is proper.  The retained
line on \(S\) becomes \(g^*R_U\), giving \(\OO_S(S)\simeq g^*R_U\).
Normalization and the periodicity cover are finite also in the analytic
category.  The separating blowup is projective.  A common
power of its equivariant relative ample line kills the bounded
finite stabilizer characters and descends to the coarse
quotient; relative ampleness is checked on fibers after the
finite pullback.  Thus these operations and a projective log
resolution admit relative ample metrics.  Adding a sufficiently
large pullback of the ambient Kähler form gives Kähler forms
near the compact sets in question.  Finally, power-compatible
root comparisons identify the full constructions as ambient
germs.  Smooth-functorial projective resolution of the same
ordered divisor data preserves those identifications.  This
proves all assertions.
\end{proof}

The residue assertions below are vacuous for an empty chart.  Fix a
nonempty chart of Lemma~\ref{signed:root-charts}, write \(Z=Z_U\), and
choose a projective log resolution \(p:\widehat Z\to Z\).  Write
\begin{equation}\label{signed:resolved-divisors}
\begin{gathered}
 D_S=p^*S,\qquad H=(D_S)_{\mathrm{red}},\qquad
 C=(p^{-1}T)_{\mathrm{red}},\\
 A=\sum_{\substack{E\text{ component of }C\\E\text{ not a component of }H}}E.
\end{gathered}
\end{equation}
The reduced divisor \(H+A\) has simple normal crossings.
The divisor \(A\) may meet \(H\); it has no component in
common with \(H\).  Thus \(A\) retains the components over \(T\)
whose poles are not already counted by the reduced divisor \(H\).
Let \(p_H=p|_H:H\to S\).
The maps that will be used in the residue and lifting arguments are
\[
\begin{tikzcd}[column sep=large,row sep=large]
H \arrow[r,hook] \arrow[d,"p_H"'] &
\widehat Z \arrow[d,"p"]\\
S \arrow[r,hook] \arrow[d,"\pi|_S"'] &
Z \arrow[d,"\pi"]\\
D_U \arrow[r,hook] \arrow[d,"\phi"'] & X\\
U &
\end{tikzcd}
\qquad g=\phi\pi|_S,\qquad h=g p_H.
\]
Only the support spaces in the left column map to \(U\).

\begin{lemma}[Residue with residual poles]\label{signed:residue}
For every \(i\geq0\), the fixed adjoint isomorphism gives a
split injection
\begin{equation}\label{signed:residue-injection}
 R^ig_*\OO_S(aS)\lhook\joinrel\longrightarrow
 R^ih_*\omega_H(A).
\end{equation}
The injection and its retraction commute with the root germ
comparisons.  Under a product with a manifold they are the
relative canonical constructions; for absolute canonical
sheaves they are tensored with the canonical line of that
additional factor.
\end{lemma}

\begin{proof}
The quotient retraction follows the construction of
\cite[Proposition~10.4]{BL}; the residual divisor changes
the comparison sheaf, which we now compute.
Apply Equation~\eqref{signed:adjoint-chart} to the rational
section \(1\) of \(\OO_Z(aS)\), and denote the resulting
fixed meromorphic adjoint form again by \(\tau\).  Its pullback has at
most a logarithmic pole along \(H+A\).  This is the log
discrepancy inequality for \((Z,S+T)\); away from the inverse
boundary, the klt ambient property and integrality of the
orders remove a possible pole.  We obtain
\[
 p^*\OO_Z(aS)\longrightarrow\omega_{\widehat Z}(H+A).
\]
Residue on the reduced SNC divisor \(H\) gives a morphism
\begin{equation}\label{signed:derived-insertion}
 \OO_S(aS)[0]\longrightarrow R(p_H)_*\omega_H(A)
\end{equation}
in the derived category of \(S\).

The insertion counts a component common to \(C\) and \(H\) in the
logarithmic divisor \(H\).  The retraction compares with \(\omega_Z(T)\),
so it must retain every component over \(T\); it therefore uses the full
divisor \(C\).  We claim
\begin{equation}\label{signed:full-residual-comparison}
 Rp_*\omega_{\widehat Z}(C)=\omega_Z(T)[0].
\end{equation}
The right side is invertible, since it is
\(\omega_Z(S+T)(-S)\).  A local frame pulls back with at
most simple poles over \(T\): subtracting the effective
Cartier \(S\) from the lc boundary leaves the required lc
order inequality.  At a prime not over \(T\), the target is
klt with invertible canonical sheaf, so its integral order is
nonnegative.  This proves the inclusion of the right side in
\(p_*\omega_{\widehat Z}(C)\).  Conversely, orders at the
strict transforms give at most a simple pole on each prime of
\(T\) and none on any other prime.  Normality and reflexivity
give the opposite inclusion.

For higher direct images this can be checked locally on \(Z\).
Choose an effective Cartier divisor \(V\) whose support is
\(T\), and a sufficiently small positive rational \(\epsilon\)
so that \(C=\lceil\epsilon p^*V\rceil\).  There are only
finitely many relevant components after shrinking about a
compact fiber.  The fractional support is SNC, and
\(\epsilon p^*V\) is \(p\)-nef and \(p\)-big.  The latter
can be seen directly after shrinking the target to a Stein
open.  Choose a \(p\)-ample Cartier divisor \(H_p\).  The
coherent sheaf \(p_*\OO_{\widehat Z}(-H_p)\) has generic
rank one because \(p\) is birational.  Cartan's Theorem~A
supplies a section nonzero at the generic points, giving an
effective divisor \(B_p\sim-H_p\) on the inverse image.
Thus \(\epsilon p^*V\sim H_p+(B_p+\epsilon p^*V)\), a
relatively ample divisor plus an effective one, which is
\(p\)-big.
Relative Kawamata--Viehweg vanishing for projective analytic
morphisms with smooth source
\cite[Theorem~5.1]{FujinoAnalyticBCHM} gives
\(R^ip_*\omega_{\widehat Z}(C)=0\) for \(i>0\).  This
proves Equation~\eqref{signed:full-residual-comparison}.

Projection formula gives the same comparison after adding
\(D_S=p^*S\).  The quotient sequence therefore identifies
\[
\begin{aligned}
 R(p_D)_*\left(
 \frac{\omega_{\widehat Z}(D_S+C)}{\omega_{\widehat Z}(C)}
 \right)
 &\simeq \omega_Z(S+T)|_S[0]\\
 &\simeq \OO_S(aS)[0],
\end{aligned}
\]
where \(p_D:D_S\to S\) and the quotient is a sheaf on the
possibly nonreduced divisor \(D_S\).  To justify the displayed
derived statement, first push it by the exact closed immersion
\(S\hookrightarrow Z\).  There it is the quotient of the two
acyclic comparisons.  Closed pushforward detects cohomology
sheaves, and canonical truncation yields the statement on \(S\).

There is a quotient morphism
\begin{equation}\label{signed:quotient-map}
\begin{aligned}
 \omega_H(A)&=
 \frac{\omega_{\widehat Z}(H+A)}{\omega_{\widehat Z}(A)}
 \\
 &\longrightarrow
 \frac{\omega_{\widehat Z}(D_S+C)}{\omega_{\widehat Z}(C)}.
\end{aligned}
\end{equation}
Indeed \(H+A\leq D_S+C\) and \(A\leq C\).  This morphism
need not be injective when \(C\) and \(H\) share a component.
Compose its derived pushforward with the preceding isomorphism.
The result retracts Equation~\eqref{signed:derived-insertion}:
at each generic point of \(S\), its composite is the identity
under the fixed residue convention.  The composite is an
endomorphism of the invertible sheaf \(\OO_S(aS)\) in degree
zero; since \(S\) is reduced, agreement at every generic point
implies agreement everywhere.  Applying \(Rg_*\) proves
Equation~\eqref{signed:residue-injection}.  All maps use the
fixed adjoint isomorphism, divisor ideals, and residue, so they
commute with the germ comparisons.  Relative canonical forms
under products, followed by wedging with the canonical frame
of the new factor, give the last assertion.
\end{proof}

\subsection{The localized Hodge argument}

The target in Equation~\eqref{signed:residue-injection} contains
the additional poles \(A\).  The following version of the filtered
SNC calculation incorporates them.  In its application, we embed \(H\)
by the graph of \(h\) in \(\widehat Z\times U\) and use the projection to
\(U\) as the ambient map.  This projection is submersive, while its
restriction to the graph is the proper map \(h\); it requires no extension
of \(h\) to \(\widehat Z\).  We will use the same graph construction after
a change of parameter.  The parameter is projective only in the last
assertion below; the closed-stratum maps are proper Kähler maps.

\begin{proposition}\label{signed:localized-hodge}
Let \(q:\mathscr V\to Y\) be a holomorphic map of complex
manifolds, submersive near a reduced closed analytic subspace
\(i:H\hookrightarrow\mathscr V\) of pure dimension \(d\) and
codimension \(c_0>0\).  Suppose \(h=q|_H\) is proper.  Assume
that locally \(H\) is an SNC divisor in a smooth submanifold of
codimension \(c_0-1\), and that a reduced divisor \(A\) is
simultaneously SNC and transverse to those support equations.
Assume that the components \(H_i\) and \(A_j\) near \(H\) have
finite global smooth indexing, and that every closed stratum
\(H_I\cap A_J\), with \(I\ne\varnothing\), is smooth and proper
over \(Y\) and carries a global relative Kähler form.  Empty or
disconnected strata are allowed.
Here \(H_I=\bigcap_{i\in I}H_i\) and
\(A_J=\bigcap_{j\in J}A_j\), with \(A_\varnothing=\mathscr V\).

Use right \(\mathscr D\)-modules and finite-pole local cohomology,
and put
\[
 \mathcal K_A=\mathcal H^{c_0}_{[H]}(\omega_{\mathscr V})(*A).
\]
Its order filtration is generated from
\(F_0\mathcal K_A=i_*\omega_H(A)\), with \(F_p=0\) for \(p<0\).
For \(M^j=\mathcal H^j q_+\mathcal K_A\), the filtered direct
image is strict at every level, and \(M^j\) has a finite
filtration by submodules, strict for \(F\), whose quotients are
polarizable real pure Hodge modules.  In particular,
\begin{equation}\label{signed:lowest-step}
 F_0M^j=R^jh_*\omega_H(A)\lhook\joinrel\longrightarrow M^j,
 \qquad F_pM^j=0\quad(p<0).
\end{equation}
Define the symbol morphism
\[
 \sigma:F_0M^j\otimes T_Y\longrightarrow\gr^F_1M^j,
 \qquad \sigma(m\otimes\xi)=[m\xi].
\]
If \(Y\) is smooth projective, then for every ample line \(N\) on \(Y\),
\begin{equation}\label{signed:symbol-vanishing}
 \Hom_{\OO_Y}(N,\Ker\sigma)=0.
\end{equation}
This includes any coherent torsion in the kernel.
\end{proposition}

\begin{proof}
We verify the residual localization in the proof of
\cite[Proposition~11.1]{BL}.  We first identify the lowest order step and
a filtration by the number of branch poles.  Its graded pieces will be
dual constant modules on smooth strata, to which proper Kähler direct
image applies.  The resulting strictness gives the injection at the lowest
filtration step,
and the graded de Rham complex on a projective base gives the symbol
vanishing.

At a point of \(H\), choose
simultaneous coordinates
\[
 (z_1,\ldots,z_{c_0-1},x_1,\ldots,x_t,y_1,\ldots,y_s,\xi),
 \quad
 \mathscr I_H=(z_1,\ldots,z_{c_0-1},x_1\cdots x_t),
 \quad A=(y_1\cdots y_s=0).
\]
Let \(\eta\) be the coordinate volume form.  The localization
Čech complex for the displayed regular support equations has
cohomology only in degree \(c_0\).  Localizing this cohomology
also in the \(y\)'s gives \(\mathcal K_A\).  Successive finite
Taylor divisions give its unique additive polar normal forms:
finite sums of
\begin{equation}\label{signed:polar-form}
 f(x_{I^c},y_{J^c},\xi)
 \prod_{\nu=1}^{c_0-1}z_\nu^{-u_\nu}
 \prod_{i\in I}x_i^{-v_i}
 \prod_{j\in J}y_j^{-w_j}\eta,
 \quad I\ne\varnothing,
\end{equation}
where every displayed order is positive and \(J\) is arbitrary,
including empty.  The coefficient is independent of exactly
the variables occurring negatively.  The normal form is an
additive statement, not an \(\OO_{\mathscr V}\)-linear
decomposition.

The simple fractions, with each displayed order equal to one,
are exactly the image of \(\omega_H(A)\).  More explicitly,
their unreduced representatives are
\[
 \frac{f\eta}
 {z_1\cdots z_{c_0-1}(x_1\cdots x_t)(y_1\cdots y_s)}.
\]
The regular-equation residue identifies their numerators modulo
\(\mathscr I_H\) with the absolute dualizing sheaf of \(H\),
twisted by \(A\).  This map is injective.  If the fraction has
no remaining \(x\)-pole, Taylor division says that its
numerator modulo the \(z\)'s is divisible by \(x_1\cdots x_t\).
Localization in the \(y\)'s does not enlarge this kernel,
because each \(y_j\) is a nonzerodivisor modulo
\(\mathscr I_H\).  This also proves the identification under
changes of coordinates, including the determinant of the
normal conormal frame.

Give a term of Equation~\eqref{signed:polar-form} the excess
\[
 e=\sum_\nu(u_\nu-1)+\sum_{i\in I}(v_i-1)
                   +\sum_{j\in J}(w_j-1).
\]
The right canonical action of a coordinate derivative is minus
differentiation of the coefficient of \(\eta\).  A derivative
in a denominator variable increases its order by one with a
nonzero scalar.  Because the coefficient in the normal form
is independent of those variables, every term of excess
\(e\) is obtained from a simple term by \(e\) such derivatives.
Conversely, order \(p\) differentiations produce excess at
most \(p\).  Thus \(F_p\mathcal K_A\) consists exactly of the
finite sums of excess at most \(p\).  This is the good order
filtration generated from \(i_*\omega_H(A)\).

There is also an increasing filtration \(W\) by the total
number of branch poles.  Intrinsically, its step of index
\(-d+k-1\) is the sum of the images of support modules for
subunions of the \(H_i\), localized along sublists of the
\(A_j\), for which the total number of selected branches is
at most \(k\).  In the normal form it imposes
\(|I|+|J|\leq k\).  The induced double filtration therefore
has
\begin{equation}\label{signed:localized-graded}
 \gr^W_{-d+k-1}(\mathcal K_A,F)
 \simeq
 \bigoplus_{\substack{|I|+|J|=k\\ I\ne\varnothing}}
 (i_{I,J})_+(\omega_{H_I\cap A_J},F^{(0)}).
\end{equation}
Here \(F^{(0)}_p\omega_{H_I\cap A_J}=0\) for \(p<0\) and is
\(\omega_{H_I\cap A_J}\) for \(p\geq0\).
Here \(A_\varnothing=\mathscr V\), and empty intersections
contribute zero.  To check this formula, retain the polar type
using exactly \(I,J\).  Its normal denominators are the
\(c_0-1\) equations \(z_\nu\), the \(|I|\) selected \(x\)'s,
and the \(|J|\) selected \(y\)'s.  They are the regular
equations of the smooth stratum \(H_I\cap A_J\), whose
dimension is \(d-k+1\).  Right closed transfer has no
codimension shift in \(F\), and the excess counts exactly its
normal derivative orders.  Mayer--Vietoris boundary maps for
the \(x\)-subunions and the localization boundary for each
selected \(y\) identify these quotients intrinsically with
the support module of that intersection.  For the latter
assertion, localization modulo the nonlocalized module in one
transverse \(y\) direction is its first local cohomology in
that direction.  A fixed order of the component labels fixes
the residue signs, so the identifications glue.  The smooth-
support graded modules are regular holonomic, and their finite
extension \(\mathcal K_A\) is regular holonomic as well.

The filtration \(W\) has a real realization.  Tensor the
finite-pole de Rham comparisons for the regular support
coordinates with the open localization comparison for the
\(y\) coordinates, whose one-coordinate factor is the
cohomology of a punctured disk.  These comparisons are
natural for all sublists.  The resulting real support and
localization complexes complexify to the de Rham complexes
just computed.  The latter are perverse by regular
holonomicity; exact and faithful complexification gives the
same assertion for the real complexes.  Their real perverse
images then give \(W\).
On a graded stratum of dimension \(d-k+1\), the real object
is its dual constant object
\(\R^H[d-k+1](d-k+1)\).  Its first right filtration step is
zero and its weight is \(-d+k-1\), in agreement with
Equation~\eqref{signed:localized-graded}.  Residue signs and
constant normalizations preserve its real polarization.  This
is the real comparison in \cite[Section~11.2]{BL}, with the
punctured-disk factors supplying the additional localization
boundaries.

The proper-support check also survives localization.  If a
holomorphic germ \(f\) vanishes on reduced \(H\), then
\begin{equation}\label{signed:support-condition}
 fF_p\mathcal K_A\subset F_{p-1}\mathcal K_A.
\end{equation}
For a generator \(z_\nu\) of \(\mathscr I_H\), multiplication
reduces its pole order or kills the class.  Multiplication by
\(x_1\cdots x_t\) also lowers the excess or kills the class,
because the polar type always has \(I\ne\varnothing\).
The \(y\)-poles do not affect this argument.  In particular
every graded stratum stays in the proper support \(H\), and
properness is needed only there.

We have now identified the order filtration, its real pure stratum pieces,
and the proper support of each level.  We next pass these data to the base.
Apply Saito's constant-source Kähler direct-image theorem
\cite[Theorem~1 and Remark~1.2]{SaitoKahler} to each smooth
closed stratum in Equation~\eqref{signed:localized-graded},
using its stipulated relative Kähler form.  The resulting
degree-\(j\) direct image of the step indexed by \(\lambda\)
is strict and polarizable real pure of weight \(\lambda+j\).
In the submersion neighborhood, the level-\(F_p\) relative
right de Rham complex has term
\[
 F_{p-e}\mathcal K_A\otimes\bigwedge^e T_{\mathscr V/Y}
 \quad\text{in degree }-e.
\]
The simultaneous excess and branch bounds show that passing
to a \(W\)-quotient commutes with taking each such level.

To prove strictness, use the finite \(W\) spectral sequence
\[
 E_1^{-\lambda,j+\lambda}
   =\mathcal H^j q_+\gr^W_\lambda\mathcal K_A
 \quad\Longrightarrow\quad \mathcal H^j q_+\mathcal K_A.
\]
Its first-page terms and their level filtrations are strict
pure objects by the preceding application.  The differentials
are real and filtration preserving, because the support and
localization comparisons are real and natural.  The first
differential preserves weight, hence is a strict morphism of
pure Hodge modules; its kernels and cokernels remain pure.
A higher differential sends \((\lambda,j)\) to
\((\lambda-k,j+1)\) for \(k\geq2\) and strictly lowers
weight.  Such a real filtered map is zero.  On distinct strict
supports this follows from strict support; on the same dense
smooth support a real map preserves both Hodge filtrations,
and a source vector of type \((p,q)\) would have image in
\(F^p\cap\overline F^q\) of a pure object of smaller weight,
which is zero.  Both the unfiltered and each level spectral
sequence therefore degenerate at the second page.  Their
comparison is injective there.  A least-\(W\)-step argument
on the finite abutment proves injection at every \(F\) level
and strictness of its induced \(W\) filtration.  This is the
strict-direct-image argument of \cite[Proposition~11.1]{BL},
now justified for every graded term with \(A\)-poles.

For \(p<0\), the relative complex is zero.  For \(p=0\), its
only term is \(F_0\mathcal K_A=i_*\omega_H(A)\) in degree
zero.  The level injection gives
Equation~\eqref{signed:lowest-step}.

Suppose finally that \(Y\) is smooth projective.  Each strict-
support pure quotient just obtained is of exactly the type in
\cite[Lemma~11.2]{BL}: a polarizable real pure direct-image
piece arising from a dual constant on a smooth Kähler source.
That lemma compares its actual right filtration with the
Sabbah--Schnell pure component extending the same generic
polarized real variation.  Its proof uses uniqueness of the
intermediate extension and the reconstruction of \(F\) from
the canonical \(V\)-filtration, so it depends only on this
pure piece, not on the divisor presentation of its source.
It gives the negative-ample vanishing
\[
 \mathbb H^u\bigl(Y,N^{-1}\otimes\gr^F_p\mathrm{DR}_Y Q\bigr)=0
 \quad(u<0)
\]
for each pure quotient \(Q\), by
\cite[Theorem~16.3.10]{SabbahSchnell}.  The finite strict
filtration extends this vanishing to \(M^j\).  Since its
negative \(F\) levels vanish, its degree-one graded right
de Rham complex is exactly
\[
 \gr^F_1\mathrm{DR}_Y M^j=
 \left[F_0M^j\otimes T_Y\xrightarrow{\sigma}\gr^F_1M^j\right]
 \quad\text{in degrees }-1,0.
\]
Its negative hypercohomology after tensoring by \(N^{-1}\)
is \(H^0(Y,N^{-1}\otimes\Ker\sigma)\).  It vanishes, which
is Equation~\eqref{signed:symbol-vanishing}.  This calculation
does not assume that the kernel is locally free.
\end{proof}

\subsection{Lifting every finite boundary neighborhood}

Return to the charts and resolution of
Lemma~\ref{signed:root-charts} and
Equation~\eqref{signed:resolved-divisors}.  On an empty chart the sheaves
below are zero.  On a nonempty chart put \(I=\OO_Z(-S)\), an invertible
ideal.  For every integer \(j\)
and positive integer \(k\), put
\[
 \mathcal A_j=I^j/I^{j+1}=\OO_S(-jS),\qquad
 \mathcal T_{j,k}=I^j/I^{j+k}.
\]
These are sheaves on the underlying topological space of \(S\).
Pushforward by \(g\) has that meaning for \(\mathcal T_{j,k}\);
it does not require a morphism from a thickening to \(U\).
A local frame of \(R_U^{-1}\) gives a Laurent graded frame
\(u^j\) of \(\mathcal A_j\) for every \(j\).  It is a frame
on the associated graded, not an extension of a frame to \(Z\).

\begin{proposition}[All finite lifting orders]\label{signed:all-orders}
For every root chart, every \(j\in\Z\), and every \(k\geq1\),
the map of sheaves of complex vector spaces on \(U\)
\begin{equation}\label{signed:all-order-surjection}
 g_*\mathcal T_{j,k+1}\longrightarrow g_*\mathcal T_{j,k}
\end{equation}
is surjective.  The assertion holds on every parameter germ,
including germs supported at special parameters.
\end{proposition}

\begin{proof}
The assertion is immediate for an empty chart.  We induct on \(k\),
simultaneously for all integer \(j\) and all nonempty charts.  The layer
algebra and its connecting maps are the
ones in \cite[Lemma~12.1]{BL}; we recall their construction to
specify exactly which obstruction is to be killed.

Assume all smaller orders.  The connecting map for
\[
 0\longrightarrow\mathcal A_{j+k}
 \longrightarrow\mathcal T_{j,k+1}
 \longrightarrow\mathcal T_{j,k}\longrightarrow0
\]
factors uniquely through a map
\begin{equation}\label{signed:obstruction}
 \delta_k:g_*\mathcal A_j\longrightarrow R^1g_*\mathcal A_{j+k}.
\end{equation}
Indeed the smaller orders make
\(g_*\mathcal T_{j,k}\to g_*\mathcal A_j\) surjective.  For
\(k>1\), the kernel comes from \(g_*\mathcal T_{j+1,k-1}\),
by the leading-layer sequence.  The order \(k-1\) in degree
\(j+1\) lifts this kernel into \(g_*\mathcal T_{j,k+1}\), so
the current connecting map kills it.  For \(k=1\) the leading
map is the identity.  Vanishing of all maps in
Equation~\eqref{signed:obstruction} is equivalent to the
surjectivity at the current order.

The maps \(\delta_k\) form a degree-\(k\) derivation on the
Laurent graded algebra \(\bigoplus_jg_*\mathcal A_j\), with
values in its first-cohomology module.  On a parameter germ, choose a lift
of a leading section \(x\) to \(g_*\mathcal T_{j,k}\), and represent that
lift locally by \(x_\alpha\in I^j\).  Then
\(x_\beta-x_\alpha\in I^{j+k}\).
Their differences modulo \(I^{j+k+1}\) represent
\(\delta_k(x)\).  For another degree-\(l\) section \(z\),
the product difference is
\[
 x_\beta z_\beta-x_\alpha z_\alpha
 =x_\alpha(z_\beta-z_\alpha)+z_\alpha(x_\beta-x_\alpha)
 +(x_\beta-x_\alpha)(z_\beta-z_\alpha).
\]
The last term lies in \(I^{j+l+2k}\subset I^{j+l+k+1}\).
This proves the derivation rule.  If \(t_1,\ldots,t_b\) are
coordinates on \(U\), choose simultaneous representatives
\(G_{i,\alpha}\) of their lifts to \(g_*\mathcal T_{0,k}\); their
differences lie in \(I^k\).  Taylor's
formula modulo their squared differences gives, for any
holomorphic \(F\),
\begin{equation}\label{signed:chain-rule}
 \delta_k(F(t))=\sum_i F_{t_i}(t)\delta_k(t_i).
\end{equation}
Both formulas hold on germs, without removing torsion in the
target, and commute with the prescribed germ comparisons.

Put \(m_k=a+k>0\).  An order-\(k\) obstruction raises degree by \(k\),
whereas the residue injection accepts degree \(-a\).  We therefore test
leading sections \(x\) of degree \(-m_k=-(a+k)\).  Apply the split residue
injection to define
\[
 c(x)=\operatorname{res}_*\delta_k(x),\qquad
 e_i(x)=\operatorname{res}_*\bigl(x\delta_k(t_i)\bigr)
 \quad\text{in }R^1h_*\omega_H(A).
\]
Both arguments have degree \(-a\), the degree of
\(\OO_S(aS)\) in Equation~\eqref{signed:residue-injection}.
The next construction places these classes in a filtered direct-image
module.  For the identity graph, the calculation will give the undivided
relation \(c(x)+\sum_i e_i(x)\partial_{t_i}=0\).  Taking its degree-one
symbol gives \(\sigma(\sum_i e_i(x)\otimes\partial_{t_i})=0\).
We will globalize that
symbol on a projective coordinate-power cover, where the boundary root is
an ample line.  The adjugate formula expresses the symbol across the
ramification of this cover without dividing by its Jacobian determinant.
The symbol vanishing will then kill the obstructions of the parameter coordinates
on every stalk; the undivided relation will kill \(\delta_k(x)\), and the
derivation rule will recover the other degrees.

Let \(\psi:U'\to U\) be a holomorphic map between coordinate
opens of dimension \(b\), with coordinates \(w\) on \(U'\),
transverse to every smooth closed stratum \(H_I\cap A_J\).
The identity map is allowed.  In \(\widehat Z\times U'\),
the graph support \(H^\flat=H\times_U U'\) is an SNC divisor
in a smooth complete intersection of codimension \(b\).
The residual divisor remains transverse.  After shrinking about a
compact fiber, the finite-component Lemma~10.3 of \cite{BL}
gives finite smooth indexing of the graph strata, and their
relative Kähler forms restrict from the product neighborhood.
Complete-intersection adjunction gives the canonical factor
\[
 \omega_{U'}\otimes\psi^*\omega_U^{-1},
\]
whose coordinate frame we denote by \(dw/dt\).  This is a
ratio of canonical frames, with no inverse Jacobian.  Pull
the Čech representatives of \(c(x),e_i(x)\) to this graph
and tensor by that frame.  Proposition~\ref{signed:localized-hodge}
places the resulting classes \(\widetilde c,\widetilde e_i\)
in \(F_0M^\flat\), where \(M^\flat\) is the degree-one
direct-image module of the localized graph support.

Write \(J=(\partial\psi_i/\partial w_j)\) and
\(J^\#=\operatorname{adj}(J)\).  Because the adjugate is polynomial in
\(J\), the following identity remains meaningful where \(J\) is singular:
\begin{equation}\label{signed:adjugate}
 \widetilde c\det J+
 \sum_{i,j}(\widetilde e_i\partial_{w_j})J^\#_{ji}=0
 \quad\text{in }M^\flat.
\end{equation}
We verify the change from \cite[Lemma~12.2]{BL} explicitly.
Choose simultaneous representatives of the lifts to length \(k\)
\[
 x_\alpha\in I^{-a-k},\quad x_\beta-x_\alpha\in I^{-a},
 \qquad G_{i,\beta}-G_{i,\alpha}\in I^k.
\]
After pulling to \(\widehat Z\), put
\(\eta_\alpha=x_\alpha\tau\), using the meromorphic form
defined by Equation~\eqref{signed:adjoint-chart}, and put
\(Q_{i,\alpha}=\psi_i(w)-G_{i,\alpha}\),
\(P_\alpha=\prod_iQ_{i,\alpha}\).  On an overlap write
\(\Delta_i=G_{i,\beta}-G_{i,\alpha}\).  The insertion and
\(D_S\geq H\) give the precise pole bounds
\begin{align}
 \eta_\alpha&\in\omega_{\widehat Z}(H+A+kD_S),\nonumber\\
 \eta_\beta-\eta_\alpha,\quad \eta_\alpha\Delta_i
   &\in\omega_{\widehat Z}(H+A),\label{signed:pole-bounds}\\
 \eta_\alpha\Delta_i\Delta_l,\quad
 (\eta_\beta-\eta_\alpha)\Delta_i
   &\in\omega_{\widehat Z}(H+A-kD_S)
     \subset\omega_{\widehat Z}(A).\nonumber
\end{align}
Thus every quadratic or cross difference has no \(H\)-pole,
although it may retain a pole on \(A\).  Such a term is zero
in the support module localized along \(A\).

In that localized graph support module consider the finite
generalized fractions
\[
 S_\alpha=\left[\frac{\eta_\alpha\wedge dw}{P_\alpha}\right].
\]
The \(H\)-pole is already in the numerator.  Each polar class
is killed by a power of a reduced equation of \(H\), and each
\(\Delta_i\) is a multiple of that equation.
Changing \(Q_{i,\alpha}\) to \(Q_{i,\beta}\) therefore has a
finite geometric expansion on each class.  The last line of
Equation~\eqref{signed:pole-bounds} kills all terms beyond the
linear terms, and gives the exact equality
\[
 S_\beta-S_\alpha=
 \left[\frac{(\eta_\beta-\eta_\alpha)\wedge dw}{P_\alpha}\right]
 +\sum_i\left[
 \frac{\eta_\alpha\Delta_i\wedge dw}
 {Q_{i,\alpha}P_\alpha}\right].
\]
Let \(C_\bullet\) be the first cochain.  Let \(E_{i,\bullet}\)
be the cochain with the numerator of the \(i\)-th summand and
only denominator \(P_\alpha\), and let \(E_{i,\bullet}^{(l)}\)
have its additional denominator \(Q_{l,\alpha}\).  The simple
cochains \(C_\bullet,E_{i,\bullet}\) are cocycles representing
\(\widetilde c,\widetilde e_i\); their possible triple-overlap
errors are precisely the cross terms killed above.  The normal
residue frame in this identification is \(dw/dt\).

The numerator of \(E_{i,\bullet}\), including its \(A\)-poles,
is independent of \(w\).  Right differentiation therefore gives
\[
 E_{i,\bullet}\partial_{w_j}=\sum_lJ_{lj}E_{i,\bullet}^{(l)}.
\]
The preceding fraction equality is
\(\check dS_\bullet=C_\bullet+\sum_iE_{i,\bullet}^{(i)}\).
Multiplying by \(\det J\) and using \(JJ^\#=(\det J)\id\)
gives at the cochain level
\[
 \check d(S_\bullet\det J)=C_\bullet\det J+
 \sum_{i,j}(E_{i,\bullet}\partial_{w_j})J^\#_{ji}.
\]
These cochains are in the end term of the relative right de
Rham complex, so a Čech coboundary there is a total
coboundary.  Passing to the direct image proves
Equation~\eqref{signed:adjugate}, with \(J^\#_{ji}\) placed
after the right derivative as displayed.  Passing to
\(\gr^F_1\) yields
\begin{equation}\label{signed:local-symbol}
 \sigma\left(\sum_{j,i}J^\#_{ji}\widetilde e_i
                         \otimes\partial_{w_j}\right)=0.
\end{equation}
For \(\psi=\id\) the undivided identity is
\begin{equation}\label{signed:undivided}
 c(x)+\sum_i e_i(x)\partial_{t_i}=0.
\end{equation}
For \(b=0\) there are no graph equations, and the same
undivided cochain computation gives \(c(x)=0\).

We now prove that the obstructions of the parameter coordinates vanish on every
stalk, including an obstruction supported only at a special parameter.
Suppose \(b>0\), and fix \(t_*\in P\).  We will use one compact graph over
a projective parameter space, obtained by a parameter change unbranched
above \(t_*\).  The vanishing of
\(\Hom\) in Equation~\eqref{signed:symbol-vanishing} applies to the entire
coherent symbol kernel, and the local isomorphism
above \(t_*\) will detect vanishing on the full germ there.  The
construction of \cite[Proposition~10.5]{BL} gives a coordinate-
power map, here of exponent \(\ell\),
\[
\begin{gathered}
 \psi:P'=\PP^b\longrightarrow P,\qquad
 [w_0:\cdots:w_b]\longmapsto[w_0^\ell:\cdots:w_b^\ell],\\
 R=\OO_{P'}(1),\qquad R^\ell\simeq\psi^*\OO_P(1),
\end{gathered}
\]
in suitably chosen coordinates, unbranched over \(t_*\).
We spell out why its geometry also accommodates \(A\).
Choose finitely many relatively compact parameter opens
covering \(\phi(D)\), with the chart resolutions defined on
slightly larger opens.  Empty charts contribute no strata.  Properness of
the nonempty supports leaves
only finitely many closed smooth strata \(H_I\cap A_J\) meeting
the corresponding compact sets.  A general tuple of coordinate
hyperplanes is transverse to the maps of every such stratum,
for every subtuple, and avoids \(t_*\).  This follows by Sard's
theorem applied to the incidence with variable hyperplanes;
varying each hyperplane supplies its normal direction.  In
these coordinates \(d\psi\) has image the tangent space to
the intersection of its vanishing coordinate hyperplanes.
The transversality condition is consequently
\[
 dh(T_z(H_I\cap A_J))+d\psi(T_wP')=T_{h(z)}P
 \quad\text{whenever }h(z)=\psi(w).
\]
It gives the required simultaneous regular graph equations
and SNC transverse residual divisors.

To obtain a single compact graph support, first form the compact
graph \(\Xi=D\times_P P'\subset X\times P'\).  On it the
pullback of \(N_0\) has the specified root \(R\).  Lemma~10.1
of \cite{BL} extends that root to a neighborhood of \(\Xi\).
Perform the full construction of Lemma~\ref{signed:root-charts}
on this neighborhood, including the relative quotient and its ordinary
pullback, and use the canonical sheaf relative to \(P'\) in the periodicity
algebra; the absolute canonical factor is added only when
taking graph residues.  Resolve the same ordered total divisor data before imposing
the graph equations.  Near a compact graph slice, the root
comparison lemma identifies this construction with the ordinary
product of a local chart with a parameter open, as an ambient
germ preserving the ideal and the adjoint form.  Uniqueness
of the comparison as a germ and properness of the graph allow
a base shrink on which the identification holds on the entire
slice.  Normalization thus occurs before ramified graph base
change.  Smooth-functorial resolution preserves the product
identification.

The graph cut \(H'\) in this resolved ambient space is compact,
has pure dimension \(n-1\), and is locally an SNC divisor in
a smooth complete intersection of codimension \(b\).  Its
ambient projection to \(P'\) is submersive near \(H'\).
Let \(A'\) be the reduced ambient resolved residual divisor, restricted
to a neighborhood of \(H'\).  Under each product-germ identification
it is \(A\times U'\); its restriction to the graph complete intersection
is the transverse residual cut.  The chosen transversality makes \(A'\)
simultaneously SNC with the graph support equations.
All closed strata \(H'_I\cap A'_J\) are compact and smooth;
splitting their connected components leaves finitely many
indices.  The relative metric construction in
Lemma~\ref{signed:root-charts} gives one Kähler form near
\(H'\), and hence the required forms on every stratum.  Thus
Proposition~\ref{signed:localized-hodge} applies to this single
global graph.  Write \(M'\) for its degree-one module.

Take \(x=u^{-m_k}\).  On a graph chart choose a frame
\(\vartheta\) of \(R\) compatible with the downstairs frame
of \(\OO_P(1)\).  The local rule
\begin{equation}\label{signed:global-symbol}
 \vartheta^{m_k}\longmapsto
 \sum_{j,i}J^\#_{ji}\widetilde e_i(u^{-m_k})
                    \otimes\partial_{w_j}
\end{equation}
defines a global morphism
\(R^{m_k}\to F_0M'\otimes T_{P'}\).  Here is the transition
check, including the ramification locus.  For changes of
coordinates let \(A_t=\partial t'/\partial t\) and
\(B_w=\partial w'/\partial w\), and let the root frame change
by \(\vartheta'=\lambda\vartheta\).  Locally \(\lambda\)
descends from a holomorphic unit downstairs: its \(\ell\)-th
power does, and after choosing a local \(\ell\)-th root of
that unit the remaining ratio is a locally constant element
of \(\mu_\ell\).  Root germ compatibility and
Equation~\eqref{signed:chain-rule} transform the downstairs
column of \(e_i\)'s by \(\lambda^{m_k}A_t\).  No derivative
of \(\lambda\) appears, since \(x\) multiplies outside
\(\delta_k(t_i)\).  The absolute canonical factor changes
by \(\det B_w/\det A_t\).  Hence
\[
\begin{aligned}
 \widetilde e'&=\frac{\det B_w}{\det A_t}\lambda^{m_k}A_t\widetilde e,
 &J'&=A_tJB_w^{-1},\\
 (J')^\#\widetilde e'&=\lambda^{m_k}B_wJ^\#\widetilde e.
\end{aligned}
\]
The adjugate equality is polynomial in \(J\), so holds also
when \(J\) is singular.  The last formula, together with the
inverse change of the vector basis, is exactly the transition
of Equation~\eqref{signed:global-symbol}.  The graph classes
here are natural Čech pullbacks through ambient germs; no
ramified base-change isomorphism for \(R^1h_*\) is used.

Equation~\eqref{signed:local-symbol} puts the image of this
morphism in the first-symbol kernel.  Since
\(R^{m_k}=\OO_{\PP^b}(a+k)\) is ample,
Equation~\eqref{signed:symbol-vanishing} makes the morphism
zero.  At a point above \(t_*\), the map \(\psi\) is locally
biholomorphic and \(J^\#\) is invertible.  The ambient germ
comparison identifies the graph classes with the downstairs
classes there.  Thus \(e_i(u^{-m_k})=0\) as germs at
\(t_*\), including any class supported there.  The point was
arbitrary, so these classes vanish on every chart.

The split injection in Lemma~\ref{signed:residue} and the
invertible Laurent frame now give \(\delta_k(t_i)=0\).
For every leading section \(x\) of degree \(-m_k\), the
classes \(e_i(x)\) vanish.  Equation~\eqref{signed:undivided},
the injection in Equation~\eqref{signed:lowest-step}, and
then the split residue injection give \(\delta_k(x)=0\).
When \(b=0\), the direct identity \(c(x)=0\) and the same
two injections give this conclusion without a graph test.
In particular, the derivation rule and characteristic zero give
\[
 0=\delta_k(u^{-m_k})=-m_k u^{-m_k-1}\delta_k(u),
 \qquad \delta_k(u)=0.
\]
For any local \(F\in g_*\OO_S\), the degree-\(-m_k\) section
\(Fu^{-m_k}\) then gives
\(0=\delta_k(Fu^{-m_k})=u^{-m_k}\delta_k(F)\).
Every graded section is \(Fu^j\), so \(\delta_k\) vanishes
in every degree.  This proves the current order and completes
the induction.
\end{proof}

\subsection{Compact deformations and the signed theorem}

We now turn the infinitesimal lifting into actual compact
subspaces.  The role of all the integer layers in
Proposition~\ref{signed:all-orders} is that both the equations
of a boundary fiber and a generator of its normal direction
can be lifted compatibly.  This is the analytic counterpart of
the compact-family step in \cite[Proposition~5.1]{LI}.

\begin{lemma}\label{signed:escaping-fibers}
Assume \(0<v<n\), set \(d=n-v\), and choose a positive
component \(D_i\) with \(\dim\phi(D_i)=r\).  There is a nonempty
open subset \(D_i^\circ\subset D_i\) with the following
property.  For every \(x\in D_i^\circ\), there are a disk
\(\Delta\) about \(0\), a proper flat analytic family
\(\mathscr F\to\Delta\) of compact subspaces of pure
dimension \(d\), and a finite morphism
\[
 \mathscr F\longrightarrow X\times\Delta
\]
over \(\Delta\), such that every nonzero fiber has image
disjoint from \(D\), while \(x\) is a limit of points in
those images as the parameter tends to zero.  Its schematic
image is flat over the disk with pure \(d\)-dimensional fibers,
whose fundamental cycles form a bounded family in the cycle
space of \(X\) after shrinking the disk.
\end{lemma}

\begin{proof}
Choose a general point of the reduced image \(\phi(D_i)\),
outside the small image in
Equation~\eqref{signed:small-intersections} and the intersections
with distinct component images.  Work on one ordinary root
chart over a neighborhood \(U\) of that point.  Its finite
map \(S\to D_U\) has a component covering a nonempty open of
\(D_i\cap D_U\).  Shrink to a smooth open \(Y\) of
\(\phi(D_i)\cap U\), remove the images of components that
do not dominate it, and remove the nonflat locus of the
proper map \(g\).  The last removal is proper by analytic
generic flatness.  The relevant \(S\) is now flat over the
smooth \(r\)-fold \(Y\), and its fibers avoid \(T\).
They are compact of pure dimension \(n-1-r=d>0\).
On the selected component the locus where it is singular or
where \(g\) has rank less than \(r\) is proper.  Removing its
image, along with the preceding proper bad subsets, under the
finite map leaves a nonempty open \(D_i^\circ\subset D_i\).
Its points have preimages at smooth points of these fiber
components.

Fix \(x\in D_i^\circ\) and a point above it in the compact
fiber \(F=g^{-1}(t)\), for \(t\in Y\).  Work on the single
ordinary chart neighborhood of that fiber from
Lemma~\ref{signed:root-charts}.  Take regular coordinates
\(t_1,\ldots,t_r\) on \(Y\) centered at \(t\), extended to
local coordinates of its smooth embedding in \(P\).  Their
pullbacks form a regular sequence on \(S\) along \(F\),
by flatness.  In addition, \(S\) is Cohen--Macaulay, so
\(F\) is Cohen--Macaulay of pure dimension \(d\).

Choose once and for all a Hausdorff open neighborhood
\(W\subset Z\) of \(F\), disjoint from \(T\), on which
\(\pi\) is quasi-finite.  Set \(I=\OO_Z(-S)\).
Proposition~\ref{signed:all-orders} lets us lift the finitely
many \(t_i\)'s to compatible germs of sections of \(\OO_Z/I^q\),
\(q\geq1\), along \(F\).  It also lifts the conormal frame
\(u\in I/I^2\) to compatible germs
\(\widetilde y_q\in I/I^{q+1}\) along \(F\); write \(\widetilde y\)
for this system.  At each order choose simultaneous
representatives on a neighborhood of \(F\) contained in
\(W\).  These neighborhoods may shrink with the order.
The leading coefficient of \(\widetilde y\) is a unit
frame, so it generates \(I\) on each finite thickening.

For \(R_q=\C[s]/(s^q)\), cut the lifted \(r\) equations in
the Cartier thickening \(qS\), and map \(s\) to
\(\widetilde y\).  Denote the resulting subspace by
\(F_q\).  It is flat over the Artin analytic point
\(\Spec R_q\), with closed fiber \(F\).  Before cutting,
multiplication by the Cartier generator identifies each
successive \(s\)-layer with \(\OO_S\); this is the
flatness criterion over \(R_q\).  Quotienting by lifts of
the closed-fiber regular sequence preserves flatness, by
the local flatness criterion.  Compatibility of the lifted
equations gives compatible reductions of all \(F_q\).

Each of these is an embedded compact deformation in the
same fixed neighborhood \(W\).  To see why the shrinking
representative neighborhoods cause no difficulty, the
support of \(F_q\) is the fixed compact set \(F\).  Its
ideal on its representative neighborhood glues with the
unit ideal on \(W\setminus F\), since the deformation is
empty on their overlap away from \(F\).  This realizes it
as a closed subspace of \(W\times\Spec R_q\).  The ideal
is coherent locally on both opens, and the family is proper
over the Artin point because its support is compact.

The Douady space of compact subspaces of \(W\) represents
proper flat embedded families \cite[Section~9]{Douady1966}.
The compatible \(F_q\)'s thus define a formal arc in its
analytic germ at \([F]\).  Embed that germ in a finite
dimensional analytic coordinate space.  The arc is a formal
solution of its convergent defining equations.  Analytic
Artin approximation \cite[Theorem~1.2]{Artin1968} gives a
convergent arc agreeing modulo \(s^2\).  Pull back the
universal family and shrink its disk to obtain
\(\mathscr F\subset W\times\Delta\), proper and flat over
\(\Delta\), with closed fiber \(F\) and the prescribed
first normal displacement.

The family may be taken to have pure \(d\)-dimensional
fibers throughout the disk.  Here is a local justification.
The central fiber is Cohen--Macaulay of dimension \(d\).
Flatness over the regular one-dimensional base makes its
parameter a nonzerodivisor and gives the depth and dimension
of the total local ring as \(d+1\) along that fiber.
Cohen--Macaulayness is open for analytic local rings, and
properness allows a disk shrink for which it holds on the
whole family.  Every total component then meets the central
fiber after shrinking, and flatness makes it dominate the
disk; the dimension formula gives it dimension \(d+1\).
Quotient by the nonzerodivisor at each
parameter gives Cohen--Macaulay fibers of dimension \(d\).
The analytic dimension formula and unmixedness of these
local rings give the asserted purity.

Let \(f\) be a local equation of \(S\).  On \(\mathscr F\)
it vanishes on the central fiber, so flatness writes it as
\(f=s h\).  Agreement modulo \(s^2\) with the formal
deformation says that \(h|_F\) is a unit: this is precisely
the lifted conormal generator.  Finitely many such
neighborhoods cover the compact \(F\).  Properness of
\(\mathscr F\to\Delta\) lets us shrink the disk so that
they cover the entire family and all their \(h\)'s are
units; remove the closed image of the complement.  Thus a
nonzero fiber misses \(S\).  The neighborhood \(W\) already
misses \(T\), and the set equality in Lemma~\ref{signed:root-charts} gives
\(\pi^{-1}(|D|)\cap W=|S|\cap W\).  Hence the
nonzero fiber images miss \(D\).

The induced map \(\Pi:\mathscr F\to X\times\Delta\) is proper:
the source is proper over \(\Delta\) and the target is
Hausdorff over it, so the graph is closed and the projection
is proper.  It is quasi-finite by the choice of \(W\),
hence finite.  It therefore preserves the dimension of
every fiber component.  Every total component through a central
point dominates the disk by flatness, and its punctured part is dense.
Thus every central point, including the point chosen above \(x\),
is a limit of points of nonzero fibers of dimension \(d\).

Let \(\mathscr Z\subset X\times\Delta\) be the schematic
image of this finite map.  Its algebra is the coherent image
of \(\OO_{X\times\Delta}\to\Pi_*\OO_{\mathscr F}\).
Multiplication by a nonzero
germ from the disk is injective on \(\OO_{\mathscr F}\),
by flatness, and remains injective after the exact finite
pushforward.  It is therefore injective on the image algebra.
That algebra is torsion-free, hence flat, over each local
discrete valuation ring of the disk.  Thus \(\mathscr Z\)
is a proper flat family of subspaces of \(X\).  The finite
map to its schematic image is surjective, also on each
fiber as a map of supports.  Each \(\mathscr Z_s\) consequently
has pure dimension \(d\).  The Douady-to-cycle morphism, in the
form recorded in \cite[Section~3.3]{Fujiki1978}, makes their
fundamental cycles an analytic family in the cycle space
of \(X\).  The restriction to a smaller closed disk is
compact.  The volumes of these cycles against a Kähler
form are bounded by continuity, proving the last assertion.
\end{proof}

\begin{proof}[Proof of Theorem~\ref{signed:torsion}]
The assertion is immediate in dimension zero, and the case
\(D=0\) follows directly from the signed equivalence.  By
Lemma~\ref{signed:boundary-geometry}, the case \(v=0\)
already gives \(L\sim_{\Q}0\), and \(v=n>0\) contradicts
\(a(X)=0\).  Suppose, therefore, that \(0<v<n\).  Choose
the positive \(D_i\) in that lemma and put \(d=n-v\), so
\(0<d<n\).  We will contradict simplicity by using the
compact deformations in Lemma~\ref{signed:escaping-fibers}.

Let \(\mathcal C_d(X)\) be the Barlet space of nonzero
effective compact \(d\)-cycles on \(X\).  It is a
second-countable analytic space, hence has only countably
many irreducible components \(B_\alpha\).  For compact
Kähler \(X\) its connected components, and therefore its
irreducible components, are compact
\cite[Theorem~4.5]{Fujiki1978}.  The universal support
\(\mathcal U_\alpha\subset B_\alpha\times X\) is a compact
analytic subspace.  Take all of its full irreducible
components \(V_{\alpha\beta}\) that contain an incidence
point belonging to an integral \(d\)-cycle disjoint from
\(D\).  There are countably many of these components,
since each compact analytic \(\mathcal U_\alpha\) has
finitely many.  Their evaluation images
\[
 W_{\alpha\beta}=\operatorname{pr}_X(V_{\alpha\beta})
\]
are closed irreducible analytic subspaces of \(X\), by
proper mapping.  Each has a point outside \(D\).  We use
the full compact components here, including their fibers
at limiting cycle parameters; the incidence over only the
open set of cycles avoiding \(D\) would not have a proper
evaluation map.

These countably many closed images cover \(D_i^\circ\).
Indeed, for \(x\in D_i^\circ\), choose the family in
Lemma~\ref{signed:escaping-fibers} and a sequence
\(x_j\to x\) in nonzero fiber images.  Let \(C_j\) be an
integral component of the fundamental cycle of the schematic image fiber
containing \(x_j\), taken with coefficient one.  The cycles
\(C_j\) avoid \(D\), and their volumes are bounded by the
volumes of the total image cycles.  Bounded cycles
on compact \(X\) have relatively compact closure in
\(\mathcal C_d(X)\) \cite[Proposition~2.10]{Fujiki1978}.
The irreducible components of an analytic space are locally
finite, so this compact closure meets only finitely many
\(B_\alpha\).  After passing to a subsequence, the cycles
lie in a fixed \(B_\alpha\); after a second subsequence,
the points \((C_j,x_j)\) lie in a fixed full component
\(V_{\alpha\beta}\).  Its evaluation image is closed, so
it contains \(x\).  This proves the covering assertion.

The analytic space \(D_i^\circ\) is locally compact and
Baire.  If no \(W_{\alpha\beta}\) contained \(D_i\), each
intersection with \(D_i^\circ\) would be a proper closed
analytic subset and hence nowhere dense.  The countable
cover just proved contradicts Baire.  Consequently some
\(W_{\alpha\beta}\) contains \(D_i\).  It also contains a
point outside \(D\).  An irreducible proper analytic
subspace of the irreducible \(n\)-fold \(X\) that contains
the prime divisor \(D_i\) must equal \(D_i\), by dimension.
It follows that \(W_{\alpha\beta}=X\).

Every point of this full incidence component lies on the
support of the effective \(d\)-cycle at its parameter,
including limiting parameters.  Thus its surjective
evaluation supplies through every point of \(X\) an
irreducible compact component of such a cycle.  Each has
dimension \(d\), strictly between zero and \(n\).  This
contradicts simplicity.  The intermediate range for \(v\)
is impossible, leaving \(L\sim_{\Q}0\).  By the meaning of
rational linear equivalence for the actual line \(L\), a
positive Cartier multiple is the trivial holomorphic line.
\end{proof}

\subsection{The simple case of the induction}

We now pass from the original pair to a reduced signed boundary and a nef
model, and then subtract the added boundary from the torsion comparison.

\begin{proof}[Proof of Proposition~\ref{simple:case}]
Let \((X,B)\) and \(J=K_X+B\) be as in that proposition.
A positive-dimensional compact complex curve is projective,
so \(a(X)=0\) implies \(n\geq2\).  By
Theorem~\ref{nv:meromorphic}, a positive multiple of
\(K_X\) has a nonzero meromorphic section.  It gives an
actual rational equivalence \(K_X\sim_{\Q}G_K\) for a
signed rational divisor \(G_K\).

Choose a projective log resolution \(p_0:W\to X\) of the
boundary and the support of this divisor.  Let \(B_W\) be
the strict transform of \(B\) together with every new
exceptional prime with coefficient one.  The resolution transfer in
Proposition~\ref{bir:resolution-transfer}
gives the actual rational identity
\[
 K_W+B_W\sim_{\Q}p_0^*J+E_0,
 \qquad E_0\geq0\quad\text{exceptional over }X.
\]
The canonical pullback formula transforms \(G_K\) to a
signed representative of \(K_W\).  Enlarge the reduced
support of \(B_W\) and this representative to a reduced
SNC divisor \(D\).  Then
\begin{equation}\label{signed:enlarged-adjoint}
 J_D:=K_W+D\sim_{\Q}p_0^*J+E_0+(D-B_W),
 \qquad E_0+(D-B_W)\geq0.
\end{equation}
In particular \(J_D\) is pseudo-effective and has a
signed rational representative supported on \(D\).
Simplicity and algebraic dimension zero persist under these
birational modifications.

Apply Corollary~\ref{prog:signed-nef} to \((W,D)\) and \(J_D\).
Equation~\eqref{signed:enlarged-adjoint} supplies pseudo-effectivity, and
the signed representative is supported on the floor \(D\).
Denote the resulting ordinary dlt nef model by
\((X_{\mathrm{nef}},D_{\mathrm{nef}})\), and put
\(L_{\mathrm{nef}}=K_{X_{\mathrm{nef}}}+D_{\mathrm{nef}}\).
As a working model of the program, \(X_{\mathrm{nef}}\) is globally
strongly \(\Q\)-factorial, and the pair has the resolution property of
Definition~\ref{def:lc-strata-resolution}.  The space \(X_{\mathrm{nef}}\)
is simple and compact Kähler with \(a(X_{\mathrm{nef}})=0\);
\(D_{\mathrm{nef}}\) is reduced, \(L_{\mathrm{nef}}\) is nef, and its
actual signed representative is supported on \(D_{\mathrm{nef}}\).
Theorem~\ref{floor:generation} makes
\(L_{\mathrm{nef}}|_{D_{\mathrm{nef}}}\)
semiample on the whole reduced floor.

We check the remaining pseudo-effectivity hypothesis of
Theorem~\ref{signed:torsion}, with exactly \(c=1\).
Take a common projective resolution
\(\alpha:V\to W\), \(\beta:V\to X_{\mathrm{nef}}\), with \(V\) smooth
compact Kähler.  The manifold \(V\) is still simple, hence
not uniruled.  Ou's criterion
\cite[Theorem~1.1]{Ou2025} makes \(K_V\) pseudo-effective;
this implication uses no assumption about the absence of
canonical sections.  Since \(X_{\mathrm{nef}}\) is globally strongly
\(\Q\)-factorial, \(K_{X_{\mathrm{nef}}}\) is a rational line.  Write its
canonical comparison as an identity of actual rational lines
\[
 K_V\sim_{\Q}\beta^*K_{X_{\mathrm{nef}}}+F_K,
 \qquad F_K=F_K^+-F_K^-,
\]
where both \(F_K^+\) and \(F_K^-\) are effective and
\(\beta\)-exceptional.  Thus
\(\beta^*\{K_{X_{\mathrm{nef}}}\}+\{F_K^+\}=\{K_V\}+\{F_K^-\}\) is
pseudo-effective.  Exceptional translation in
Lemma~\ref{bir:exceptional} removes the effective
\(\beta\)-exceptional \(F_K^+\) and proves that
\(\beta^*\{K_{X_{\mathrm{nef}}}\}\) is pseudo-effective.  Since
\(L_{\mathrm{nef}}-D_{\mathrm{nef}}=K_{X_{\mathrm{nef}}}\) as an actual
rational line, this is precisely the resolution hypothesis for
\(L_{\mathrm{nef}}-D_{\mathrm{nef}}\),
with \(c=1\).  Theorem~\ref{signed:torsion} now gives
\(L_{\mathrm{nef}}\sim_{\Q}0\).

On the same common resolution, enlarging it if necessary,
the accumulated negative-step comparisons give
\begin{equation}\label{signed:program-decomposition}
 \alpha^*J_D\sim_{\Q}\beta^*L_{\mathrm{nef}}+F,
 \qquad F\geq0\quad\text{exceptional over }X_{\mathrm{nef}}.
\end{equation}
All identities here are identities of actual rational
lines.  Because \(L_{\mathrm{nef}}\) is torsion, its pullback is nef
with zero class.  Lemma~\ref{bir:exceptional} therefore
identifies \(F=N(\alpha^*J_D)\).

Put \(\mu=p_0\alpha:V\to X\) and
\(Q=\alpha^*(E_0+D-B_W)\geq0\).  Pulling up
Equation~\eqref{signed:enlarged-adjoint} gives
\(\mu^*J\sim_{\Q}\alpha^*J_D-Q\), which is
pseudo-effective.  Apply Lemma~\ref{bir:unique-current}
to Equation~\eqref{signed:program-decomposition}.  A
positive current for \(\mu^*J\), plus \([Q]\), must be
the unique current \([F]\) in the larger adjoint class.
It follows that \(Q\leq F\), and negative-part
subtraction gives the actual identity
\[
 \mu^*J\sim_{\Q}F-Q=N(\mu^*J).
\]
The divisor on the right is effective and rational.  This
is exactly Proposition~\ref{simple:case} and the required
instance of \(\Ggood_n\) with torsion positive part.
\end{proof}
